\section{Limitations and Threats to Validity}\label{sec:limits}
\textbf{No live model and no human evidence.} All data are replayed recorded evaluator payloads for 8 constructed questions and 64 constructed answers, injected faults, or one local socket server. The agents are rule-based; no LLM output was produced; follow-ups and feedback came from the hub's template fallbacks. We make no claim about coaching quality, learning, evaluator validity, or that the agents improve PPO. \textbf{Replay candidate.} Scores are fixed per answer and follow-ups reuse the root answer, so the effect of any decision on later answers cannot appear. \textbf{Dependent sessions.} The 160 XA sessions per condition reuse 64 answers; intervals are descriptive. \textbf{Metrics.} Level changes, reversals and agreement with PPO are behaviour measures, not quality measures. \textbf{Unexercised regions.} Replay has no audio, so confidence and hesitation are constant and the divergent guardrail regions are not exercised; the XA result that canonical and legacy guardrails give identical paths is specific to this input. \textbf{Post hoc additions.} XC was designed after XA and XB showed that M4 bundles four checks; its protocol and tag precede its runs, but its motivation depends on earlier results. The cluster-bootstrap intervals for XA are post hoc. With 8 recorded answer sets, even cluster intervals are coarse. \textbf{Registered but not pre-confirmed.} A9, A10 and B6 were contradicted; three metric definitions were fixed after smoke runs (disclosed in the protocol). \textbf{Agent design.} Thresholds ($\tau$, ledger parameters) were fixed in advance, not tuned, and are heuristics; the agents are not claimed to be optimal. Persistent memory is restored from a session snapshot, but a retried turn re-uses the first observation of that turn. \textbf{Coverage.} The coaching renderer is exercised on the verbal path only; the code-submission path collects only the Strategy proposal. The Strategy Agent never proposed a harder level in these experiments (no \emph{stretch} proposal), so that rule is untested here. \textbf{Trust boundary.} The agents are in-process trusted code: the authority matrix constrains what their outputs can do (schema, forbidden keys, version, arbitration), not what their code could do, so a malicious agent implementation is out of scope. \textbf{Deployment surface.} In the shipped backend the layer is enabled only by a server-side environment variable (default off; clients cannot set it), added after the main experiments (which used the constructor option). The web client has no display of agent output; the shipped-path server run used recorded evaluator payloads. After a restart the agents' ledgers, decision records and call counters are restored from the snapshot (breaker state is not). \textbf{Environment.} Runs used one Python environment outside the repository's dependency pins. The historical 407-test baseline was not reproduced as a whole: the 28 frontend tests pass, but the Python suite has 23 failures and 36 errors (Docker sandbox, model downloads, one documented known failure), the same 59 test identifiers as in a clean export of the frozen baseline commit, so a full green run was not obtained; at the audited head 421 Python tests pass (306 in the frozen export) with no new failure. \textbf{Provenance.} Protocol tags are local and carry no external timestamp; the reference list was checked through a search interface (titles, venues and abstracts) and full texts were not read; the official venue instructions could not be retrieved from the authors' environment.
